\section{Related Work}

\subsection{Reinforcement Learning for Traffic Signal Control}
RL-based signal control has progressed from single-intersection deep
Q-networks~\cite{genders2016_drl,wei2018_intellilight} through pressure-based and
cooperative multi-agent designs---PressLight~\cite{wei2019_presslight},
MPLight~\cite{chen2020_mplight}, CoLight~\cite{wei2019_colight}, and earlier
coordinated learners~\cite{vanderpol2016_coordinated}, building on the
max-pressure principle~\cite{varaiya2013_maxpressure}---to centralized-training,
decentralized-execution (CTDE) actor--critic methods such as
MAPPO~\cite{yu2022_mappo} and IPPO~\cite{dewitt2020_ippo}, with recent
IEEE~T-ITS work adding attention-based soft actor--critic
control~\cite{mao2023_masac}, incentive-driven
communication~\cite{zhou2024_micdrl}, constrained
partitioning~\cite{gu2024_partition}, intersection
heterogeneity~\cite{bie2024_heterogeneity}, and cross-city
meta-learning~\cite{jiang2024_xlight}; surveys document this
trajectory~\cite{wei2019_survey,rltsc_review2025,rltsc_survey2026}. Despite this
architectural diversity, every member of the lineage consumes privileged
simulator state---exact halting counts, per-vehicle accumulated waiting times,
and holistic network observations---that no commodity roadside sensor can
produce. This privilege dependence also makes a direct empirical comparison
ill-posed---re-implementing PressLight, MPLight, or CoLight on camera-computable
proxies would no longer faithfully represent them---so we adopt
IPPO~\cite{dewitt2020_ippo}, the strongest learned CTDE baseline that runs on
the \emph{identical} proxy observation and reward, as the learned reference, and
separately bound the advantage available to any privileged-input method through
a privileged-observation variant of our own policy (the information-restriction
cost of Section~\ref{sec:restriction_cost}).

\subsection{Deployability- and Sensing-Constrained TSC}
A growing body of work confronts this sensing gap along three threads.
The first constrains observations to legacy point detectors: RL with partial vehicle
detection~\cite{zhang2021_partialdetection} learns under binary loop-sensor
inputs, but point sensing cannot recover the spatial vehicle composition, class
mix, or continuous speed profiles that mixed-traffic control requires. The
second, vision-based RL-TSC, exploits computer vision to derive richer
features---image-based state with a learned world model~\cite{dai2022_worldmodels},
infrastructure-camera co-simulation across multiple junctions~\cite{azfar2025_cosim},
and vision-plus-language controllers~\cite{su2025_visionllm}. Critically, these
frameworks restrict the \emph{observation} space to camera-feasible queries yet
continue to train and evaluate against a \emph{privileged reward}---simulator
delay, exact waiting time, or network throughput---that is unavailable at the
edge, so the objective optimized in simulation cannot be reproduced at
deployment. The third thread, sim-to-real transfer for TSC, targets the
\emph{dynamics} gap rather than the \emph{observation} gap: prompt-based domain
transfer~\cite{da2024_promptgat} aligns transition discrepancies, mirroring the
dynamics-randomization tradition in robotics~\cite{peng2018_dynrand,zhao2020_sim2real_survey}.
Our contribution is orthogonal to all three. We impose a strict, paired
restriction of both the state \emph{and} the reward to camera-computable
proxies, and we quantify the resulting information-restriction cost together
with the policy's degradation under a sensing-noise model whose per-feature
magnitudes are fixed on a documented basis from the detection-and-tracking
literature, with conclusions resting on the shape of the degradation across a
$0$--$2\times$ magnitude sweep (Section~\ref{sec:robustness_noise}).

\subsection{Mixed and Motorcycle-Dominant Traffic Modeling}
Traffic in developing metropolises is fundamentally heterogeneous and
motorcycle-dominant: motorcycles filter laterally between lanes and swerving
lines, behavior that violates the discrete, lane-indexed queue abstraction
underlying most RL-TSC formulations. Capturing these dynamics requires
continuous-space microsimulation, such as the sublane model of
SUMO~\cite{lopez2018_sumo}, rather than single-file lane queues. Field studies of
Vietnamese cities document the magnitude of this dominance: Ngoc \emph{et al.}~\cite{ngoc2021_fivecities}
report a motorcycle mode share of $72.6\%$ in Hanoi,
rising to $\approx 80$--$89\%$ in Hai Phong, Da Nang, Ho Chi Minh City, and Can Tho,
while microscopic trajectory analyses confirm pervasive lateral filtering under
such mixes~\cite{vn_roundabout}. These shares lie an order of magnitude above
the car-centric compositions assumed in most RL-TSC studies; we therefore adopt
the Hanoi figure as a representative anchor ($\approx 73\%$ motorcycles by
vehicle count) and bracket the field range with a composition sweep
(Section~\ref{sec:generalization}). We further distinguish count share from traffic load: at a
motorcycle passenger-car-unit (PCU) equivalence of $\approx 0.29$
\cite{cao_sano_hanoi}, the $73\%$ count share corresponds to roughly one-third
of the PCU-based load, so we report count share when analyzing lateral filtering
and PCU when discussing capacity. Despite the global prevalence of such traffic,
RL-TSC evaluated under sublane dynamics remains scarce, leaving solutions for
developing-city deployment underexplored.

Table~\ref{tab:positioning} positions representative RL-TSC studies against
criteria most relevant to field deployment: whether the \emph{state} and
\emph{reward} are camera-computable, whether evaluation spans \emph{multiple
junctions} and \emph{mixed-traffic sublane} dynamics, whether the policy is
tested for \emph{robustness to sensing noise} (perturbed or degraded
observations), and whether the \emph{information-restriction cost}---the
performance gap between privileged and camera-computable signals---is quantified. To our knowledge, the camera-computable restriction of \emph{both} state and reward, evaluated under mixed-traffic sublane dynamics with an explicit restriction-cost and sensing-noise analysis, is not jointly satisfied by prior work.

\begin{table}[t]
\centering
\caption{Positioning Matrix of Representative RL-TSC Literature}
\label{tab:positioning}
\resizebox{\columnwidth}{!}{%
\begin{tabular}{lcccccc}
\toprule
\textbf{Work} & \textbf{State} & \textbf{Reward} & \textbf{Multi-} & \textbf{Mixed/} & \textbf{Noise} & \textbf{Restr.} \\
              & \textbf{Source} & \textbf{Source} & \textbf{Junc.} & \textbf{Sublane} & \textbf{Robust.} & \textbf{Cost} \\
\midrule
PressLight~\cite{wei2019_presslight}              & Privileged   & Privileged   & \checkmark & $\times$ & $\times$ & $\times$ \\
CoLight~\cite{wei2019_colight}                    & Privileged   & Privileged   & \checkmark & $\times$ & $\times$ & $\times$ \\
Partial-Det.~RL~\cite{zhang2021_partialdetection} & Partial(V2I) & Privileged   & $\times$   & $\times$ & $\times$ & $\times$ \\
Dai \emph{et al.}~\cite{dai2022_worldmodels}      & Camera-feas. & Privileged   & $\times$   & $\times$ & $\times$ & $\times$ \\
Azfar \emph{et al.}~\cite{azfar2025_cosim}        & Camera-feas. & Privileged   & \checkmark & $\times$ & \checkmark & $\times$ \\
\textbf{Ours}                                      & \textbf{Camera-feas.} & \textbf{Camera-feas.} & \textbf{\checkmark} & \textbf{\checkmark} & \textbf{\checkmark} & \textbf{\checkmark} \\
\bottomrule
\end{tabular}%
}
\end{table}